\documentclass[10pt]{article}
% Hojas resumen: un A4 por delante y por detrás, dos columnas por cara.
% Uso: \documentclass[10pt]{article} % Hojas resumen: un A4 por delante y por detrás, dos columnas por cara.
% Uso: \documentclass[10pt]{article} % Hojas resumen: un A4 por delante y por detrás, dos columnas por cara.
% Uso: \documentclass[10pt]{article} \input{estilo-resumen.tex} \bloque{analisis|algebra|azar} \cabecera{n}{Título}{Lecciones}
\usepackage[utf8]{inputenc}
\usepackage[T1]{fontenc}
\usepackage[spanish,es-nodecimaldot]{babel}
\usepackage{lmodern}
\usepackage{amsmath,amssymb}
\usepackage{graphicx}
\usepackage[a4paper,left=10mm,right=10mm,top=15mm,bottom=13mm,headheight=11pt,headsep=3mm,footskip=6mm]{geometry}
\usepackage{multicol}
\setlength{\columnsep}{6mm}
\setlength{\columnseprule}{0.3pt}
\usepackage{xcolor}
\usepackage{enumitem}
\setlist{nosep,leftmargin=1.1em,itemsep=1pt}
\usepackage{array}
\usepackage{needspace}
\usepackage{pgfplots}
\usepgfplotslibrary{fillbetween}
\pgfplotsset{compat=1.18,
  mini/.style={width=4.3cm,height=3.3cm,scale only axis=false,axis lines=middle,xtick=\empty,ytick=\empty,samples=120,
    xlabel={\scriptsize$x$},ylabel={\scriptsize$y$},title style={font=\scriptsize,yshift=-4pt},clip=true,
    every axis plot/.append style={thick,color=acento}}}
\usepackage{ragged2e}
\usepackage[most]{tcolorbox}
\usepackage{fancyhdr,lastpage}
\usepackage[hidelinks,pdfencoding=auto]{hyperref}
\hypersetup{pdflang=es-ES,pdfsubject={Matemáticas Aplicadas a las Ciencias Sociales II}}

\definecolor{analisis}{HTML}{1D4ED8}
\definecolor{algebra}{HTML}{86198F}
\definecolor{azar}{HTML}{9A3412}
\definecolor{aviso}{HTML}{B91C1C}
\colorlet{acento}{analisis}
\newcommand{\bloque}[1]{\colorlet{acento}{#1}}

\setlength{\parindent}{0pt}
\setlength{\parskip}{2pt}
\renewcommand{\arraystretch}{1.2}
\AtBeginDocument{\RaggedRight\fontsize{9.8}{11.9}\selectfont\setlength{\abovedisplayskip}{3pt}\setlength{\belowdisplayskip}{3pt}\setlength{\abovedisplayshortskip}{1pt}\setlength{\belowdisplayshortskip}{1pt}}

\newcommand{\temanum}{}\newcommand{\tematitulo}{}
\pagestyle{fancy}
\fancyhf{}
\fancyhead[L]{\footnotesize\color{acento}\textbf{Hoja resumen}\quad Tema \temanum. \tematitulo}
\fancyhead[R]{\footnotesize Matemáticas Aplicadas a las Ciencias Sociales II --- 2.\textsuperscript{o} Bachillerato}
\fancyfoot[L]{\scriptsize IES Al-Ándalus --- Juan María Gámez Ortiz, con la ayuda de Claude y de Codex --- CC BY-NC 4.0}
\fancyfoot[R]{\scriptsize Cara \thepage\ de \pageref{LastPage}}
\renewcommand{\headrulewidth}{0.4pt}

% Cabecera de la primera cara: número, título y lecciones del tema.
\newcommand{\cabecera}[3]{%
  \renewcommand{\temanum}{#1}\renewcommand{\tematitulo}{#2}%
  \hypersetup{pdftitle={Hoja resumen. Tema #1. #2}}%
  \noindent\begin{tcolorbox}[enhanced,colback=acento!8,colframe=acento,boxrule=0.6pt,arc=2pt,left=4pt,right=4pt,top=2pt,bottom=2pt]
  \begin{minipage}[c]{1.1cm}\includegraphics[width=1cm]{logo-al-andalus.png}\end{minipage}%
  \begin{minipage}[c]{\dimexpr\linewidth-1.2cm}{\large\bfseries\color{acento}Tema #1. #2}\\[1pt]{\scriptsize #3}\end{minipage}
  \end{tcolorbox}\vspace{-2pt}}

% Apartado con filete del color del bloque.
\newcommand{\seccion}[1]{\par\needspace{5\baselineskip}\vspace{3pt}{\color{acento}\bfseries\normalsize #1}\par\vspace{-1pt}{\color{acento}\hrule height 0.5pt}\vspace{2pt}}
\newcommand{\sub}[1]{\par\needspace{3\baselineskip}\vspace{1pt}\textbf{\color{acento}#1}\enspace}

% Fórmula clave, aviso de errores y ejemplo.
\tcbset{before upper=\RaggedRight}
\newtcolorbox{formula}{enhanced,colback=acento!6,colframe=acento!70,boxrule=0.5pt,arc=1.5pt,left=3pt,right=3pt,top=0pt,bottom=0pt,before skip=2pt,after skip=2pt}
\newtcolorbox{ojo}{enhanced,colback=aviso!5,colframe=aviso!60,boxrule=0.5pt,arc=1.5pt,left=3pt,right=3pt,top=1pt,bottom=1pt,before skip=3pt,after skip=2pt,
  fonttitle=\bfseries\footnotesize,coltitle=aviso,colbacktitle=aviso!5,title={Ojo},attach title to upper={\enspace}}
\newtcolorbox{ejemplo}{enhanced,colback=white,colframe=acento!45,boxrule=0.4pt,arc=1.5pt,left=3pt,right=3pt,top=1pt,bottom=1pt,before skip=3pt,after skip=2pt,
  fonttitle=\bfseries\footnotesize,coltitle=acento,colbacktitle=white,title={Ejemplo},attach title to upper={\enspace}}
\newcommand{\dx}{\,dx}
 \bloque{analisis|algebra|azar} \cabecera{n}{Título}{Lecciones}
\usepackage[utf8]{inputenc}
\usepackage[T1]{fontenc}
\usepackage[spanish,es-nodecimaldot]{babel}
\usepackage{lmodern}
\usepackage{amsmath,amssymb}
\usepackage{graphicx}
\usepackage[a4paper,left=10mm,right=10mm,top=15mm,bottom=13mm,headheight=11pt,headsep=3mm,footskip=6mm]{geometry}
\usepackage{multicol}
\setlength{\columnsep}{6mm}
\setlength{\columnseprule}{0.3pt}
\usepackage{xcolor}
\usepackage{enumitem}
\setlist{nosep,leftmargin=1.1em,itemsep=1pt}
\usepackage{array}
\usepackage{needspace}
\usepackage{pgfplots}
\usepgfplotslibrary{fillbetween}
\pgfplotsset{compat=1.18,
  mini/.style={width=4.3cm,height=3.3cm,scale only axis=false,axis lines=middle,xtick=\empty,ytick=\empty,samples=120,
    xlabel={\scriptsize$x$},ylabel={\scriptsize$y$},title style={font=\scriptsize,yshift=-4pt},clip=true,
    every axis plot/.append style={thick,color=acento}}}
\usepackage{ragged2e}
\usepackage[most]{tcolorbox}
\usepackage{fancyhdr,lastpage}
\usepackage[hidelinks,pdfencoding=auto]{hyperref}
\hypersetup{pdflang=es-ES,pdfsubject={Matemáticas Aplicadas a las Ciencias Sociales II}}

\definecolor{analisis}{HTML}{1D4ED8}
\definecolor{algebra}{HTML}{86198F}
\definecolor{azar}{HTML}{9A3412}
\definecolor{aviso}{HTML}{B91C1C}
\colorlet{acento}{analisis}
\newcommand{\bloque}[1]{\colorlet{acento}{#1}}

\setlength{\parindent}{0pt}
\setlength{\parskip}{2pt}
\renewcommand{\arraystretch}{1.2}
\AtBeginDocument{\RaggedRight\fontsize{9.8}{11.9}\selectfont\setlength{\abovedisplayskip}{3pt}\setlength{\belowdisplayskip}{3pt}\setlength{\abovedisplayshortskip}{1pt}\setlength{\belowdisplayshortskip}{1pt}}

\newcommand{\temanum}{}\newcommand{\tematitulo}{}
\pagestyle{fancy}
\fancyhf{}
\fancyhead[L]{\footnotesize\color{acento}\textbf{Hoja resumen}\quad Tema \temanum. \tematitulo}
\fancyhead[R]{\footnotesize Matemáticas Aplicadas a las Ciencias Sociales II --- 2.\textsuperscript{o} Bachillerato}
\fancyfoot[L]{\scriptsize IES Al-Ándalus --- Juan María Gámez Ortiz, con la ayuda de Claude y de Codex --- CC BY-NC 4.0}
\fancyfoot[R]{\scriptsize Cara \thepage\ de \pageref{LastPage}}
\renewcommand{\headrulewidth}{0.4pt}

% Cabecera de la primera cara: número, título y lecciones del tema.
\newcommand{\cabecera}[3]{%
  \renewcommand{\temanum}{#1}\renewcommand{\tematitulo}{#2}%
  \hypersetup{pdftitle={Hoja resumen. Tema #1. #2}}%
  \noindent\begin{tcolorbox}[enhanced,colback=acento!8,colframe=acento,boxrule=0.6pt,arc=2pt,left=4pt,right=4pt,top=2pt,bottom=2pt]
  \begin{minipage}[c]{1.1cm}\includegraphics[width=1cm]{logo-al-andalus.png}\end{minipage}%
  \begin{minipage}[c]{\dimexpr\linewidth-1.2cm}{\large\bfseries\color{acento}Tema #1. #2}\\[1pt]{\scriptsize #3}\end{minipage}
  \end{tcolorbox}\vspace{-2pt}}

% Apartado con filete del color del bloque.
\newcommand{\seccion}[1]{\par\needspace{5\baselineskip}\vspace{3pt}{\color{acento}\bfseries\normalsize #1}\par\vspace{-1pt}{\color{acento}\hrule height 0.5pt}\vspace{2pt}}
\newcommand{\sub}[1]{\par\needspace{3\baselineskip}\vspace{1pt}\textbf{\color{acento}#1}\enspace}

% Fórmula clave, aviso de errores y ejemplo.
\tcbset{before upper=\RaggedRight}
\newtcolorbox{formula}{enhanced,colback=acento!6,colframe=acento!70,boxrule=0.5pt,arc=1.5pt,left=3pt,right=3pt,top=0pt,bottom=0pt,before skip=2pt,after skip=2pt}
\newtcolorbox{ojo}{enhanced,colback=aviso!5,colframe=aviso!60,boxrule=0.5pt,arc=1.5pt,left=3pt,right=3pt,top=1pt,bottom=1pt,before skip=3pt,after skip=2pt,
  fonttitle=\bfseries\footnotesize,coltitle=aviso,colbacktitle=aviso!5,title={Ojo},attach title to upper={\enspace}}
\newtcolorbox{ejemplo}{enhanced,colback=white,colframe=acento!45,boxrule=0.4pt,arc=1.5pt,left=3pt,right=3pt,top=1pt,bottom=1pt,before skip=3pt,after skip=2pt,
  fonttitle=\bfseries\footnotesize,coltitle=acento,colbacktitle=white,title={Ejemplo},attach title to upper={\enspace}}
\newcommand{\dx}{\,dx}
 \bloque{analisis|algebra|azar} \cabecera{n}{Título}{Lecciones}
\usepackage[utf8]{inputenc}
\usepackage[T1]{fontenc}
\usepackage[spanish,es-nodecimaldot]{babel}
\usepackage{lmodern}
\usepackage{amsmath,amssymb}
\usepackage{graphicx}
\usepackage[a4paper,left=10mm,right=10mm,top=15mm,bottom=13mm,headheight=11pt,headsep=3mm,footskip=6mm]{geometry}
\usepackage{multicol}
\setlength{\columnsep}{6mm}
\setlength{\columnseprule}{0.3pt}
\usepackage{xcolor}
\usepackage{enumitem}
\setlist{nosep,leftmargin=1.1em,itemsep=1pt}
\usepackage{array}
\usepackage{needspace}
\usepackage{pgfplots}
\usepgfplotslibrary{fillbetween}
\pgfplotsset{compat=1.18,
  mini/.style={width=4.3cm,height=3.3cm,scale only axis=false,axis lines=middle,xtick=\empty,ytick=\empty,samples=120,
    xlabel={\scriptsize$x$},ylabel={\scriptsize$y$},title style={font=\scriptsize,yshift=-4pt},clip=true,
    every axis plot/.append style={thick,color=acento}}}
\usepackage{ragged2e}
\usepackage[most]{tcolorbox}
\usepackage{fancyhdr,lastpage}
\usepackage[hidelinks,pdfencoding=auto]{hyperref}
\hypersetup{pdflang=es-ES,pdfsubject={Matemáticas Aplicadas a las Ciencias Sociales II}}

\definecolor{analisis}{HTML}{1D4ED8}
\definecolor{algebra}{HTML}{86198F}
\definecolor{azar}{HTML}{9A3412}
\definecolor{aviso}{HTML}{B91C1C}
\colorlet{acento}{analisis}
\newcommand{\bloque}[1]{\colorlet{acento}{#1}}

\setlength{\parindent}{0pt}
\setlength{\parskip}{2pt}
\renewcommand{\arraystretch}{1.2}
\AtBeginDocument{\RaggedRight\fontsize{9.8}{11.9}\selectfont\setlength{\abovedisplayskip}{3pt}\setlength{\belowdisplayskip}{3pt}\setlength{\abovedisplayshortskip}{1pt}\setlength{\belowdisplayshortskip}{1pt}}

\newcommand{\temanum}{}\newcommand{\tematitulo}{}
\pagestyle{fancy}
\fancyhf{}
\fancyhead[L]{\footnotesize\color{acento}\textbf{Hoja resumen}\quad Tema \temanum. \tematitulo}
\fancyhead[R]{\footnotesize Matemáticas Aplicadas a las Ciencias Sociales II --- 2.\textsuperscript{o} Bachillerato}
\fancyfoot[L]{\scriptsize IES Al-Ándalus --- Juan María Gámez Ortiz, con la ayuda de Claude y de Codex --- CC BY-NC 4.0}
\fancyfoot[R]{\scriptsize Cara \thepage\ de \pageref{LastPage}}
\renewcommand{\headrulewidth}{0.4pt}

% Cabecera de la primera cara: número, título y lecciones del tema.
\newcommand{\cabecera}[3]{%
  \renewcommand{\temanum}{#1}\renewcommand{\tematitulo}{#2}%
  \hypersetup{pdftitle={Hoja resumen. Tema #1. #2}}%
  \noindent\begin{tcolorbox}[enhanced,colback=acento!8,colframe=acento,boxrule=0.6pt,arc=2pt,left=4pt,right=4pt,top=2pt,bottom=2pt]
  \begin{minipage}[c]{1.1cm}\includegraphics[width=1cm]{logo-al-andalus.png}\end{minipage}%
  \begin{minipage}[c]{\dimexpr\linewidth-1.2cm}{\large\bfseries\color{acento}Tema #1. #2}\\[1pt]{\scriptsize #3}\end{minipage}
  \end{tcolorbox}\vspace{-2pt}}

% Apartado con filete del color del bloque.
\newcommand{\seccion}[1]{\par\needspace{5\baselineskip}\vspace{3pt}{\color{acento}\bfseries\normalsize #1}\par\vspace{-1pt}{\color{acento}\hrule height 0.5pt}\vspace{2pt}}
\newcommand{\sub}[1]{\par\needspace{3\baselineskip}\vspace{1pt}\textbf{\color{acento}#1}\enspace}

% Fórmula clave, aviso de errores y ejemplo.
\tcbset{before upper=\RaggedRight}
\newtcolorbox{formula}{enhanced,colback=acento!6,colframe=acento!70,boxrule=0.5pt,arc=1.5pt,left=3pt,right=3pt,top=0pt,bottom=0pt,before skip=2pt,after skip=2pt}
\newtcolorbox{ojo}{enhanced,colback=aviso!5,colframe=aviso!60,boxrule=0.5pt,arc=1.5pt,left=3pt,right=3pt,top=1pt,bottom=1pt,before skip=3pt,after skip=2pt,
  fonttitle=\bfseries\footnotesize,coltitle=aviso,colbacktitle=aviso!5,title={Ojo},attach title to upper={\enspace}}
\newtcolorbox{ejemplo}{enhanced,colback=white,colframe=acento!45,boxrule=0.4pt,arc=1.5pt,left=3pt,right=3pt,top=1pt,bottom=1pt,before skip=3pt,after skip=2pt,
  fonttitle=\bfseries\footnotesize,coltitle=acento,colbacktitle=white,title={Ejemplo},attach title to upper={\enspace}}
\newcommand{\dx}{\,dx}

\bloque{azar}
\begin{document}
\cabecera{9}{Inferencia estadística}{Lecciones 9.1--9.9: estadística inferencial, muestreo y sus técnicas, distribución muestral y teorema central del límite, estimación, intervalos característicos, intervalos de confianza, tamaño de la muestra y ficha técnica.}
\begin{multicols}{2}

\seccion{1. Población y muestra}
\begin{itemize}
  \item \textbf{Población}: todos los individuos que interesan. \textbf{Muestra}: la parte que se observa.
  \item \textbf{Parámetro}: valor de la población, normalmente desconocido (\(\mu\), \(\sigma\), \(p\)).
  \item \textbf{Estadístico}: valor calculado en la muestra (\(\bar x\), \(\hat p\)); cambia de una muestra a otra.
\end{itemize}
La inferencia usa la muestra para \textbf{estimar} parámetros con un margen de error y un nivel de confianza.

\seccion{2. Técnicas de muestreo}
\textbf{Probabilísticas} (cada individuo tiene una probabilidad conocida de ser elegido):
\begin{itemize}
  \item \textbf{Aleatorio simple}: todos con la misma probabilidad (sorteo).
  \item \textbf{Sistemático}: uno de cada \(k=\frac Nn\), desde un inicio al azar.
  \item \textbf{Estratificado}: grupos homogéneos (estratos) y muestra en cada uno.
  \item \textbf{Por conglomerados}: se sortean grupos completos (aulas, barrios).
\end{itemize}
\textbf{No probabilísticas}: por cuotas, de conveniencia, bola de nieve, intencional. Pueden tener sesgo y no permiten calcular el error con estas fórmulas.
\begin{formula}
\textbf{Afijación proporcional:}
\[n_h=n\cdot\frac{N_h}{N}\]
\end{formula}
\begin{ejemplo}
Centro con 1200 alumnos: 500, 400 y 300 por nivel. Muestra de 60: \(60\cdot\frac{500}{1200}=25\), \(60\cdot\frac{400}{1200}=20\), \(60\cdot\frac{300}{1200}=15\). Si salen decimales, se redondea conservando el total.
\end{ejemplo}

\seccion{3. Distribución de la media muestral}
\begin{formula}
\[E(\bar X)=\mu,\qquad \operatorname{DE}(\bar X)=\frac{\sigma}{\sqrt n}\]
\[\bar X\sim N\Bigl(\mu,\ \frac{\sigma}{\sqrt n}\Bigr)\]
\end{formula}
Exacta si la población es normal; aproximada por el \textbf{teorema central del límite} si \(n\ge30\), sea cual sea la población.
\begin{ejemplo}
Estaturas con \(\mu=170\) cm y \(\sigma=10\) cm; muestras de \(n=25\): \(\frac{\sigma}{\sqrt n}=2\).
\[P(\bar X\ge173)=1-\Phi\Bigl(\frac{173-170}{2}\Bigr)=1-\Phi(1{,}5)=0{,}0668\]
\end{ejemplo}
\begin{ojo}
La media muestral varía mucho menos que un individuo: \(P(X\ge173)=1-\Phi(0{,}3)\approx0{,}38\), pero \(P(\bar X\ge173)\approx0{,}07\).
\end{ojo}

\seccion{4. Distribución de la proporción muestral}
\begin{formula}
\[\hat P\approx N\Bigl(p,\ \sqrt{\frac{pq}n}\Bigr),\qquad q=1-p\]
\end{formula}
Si \(n\) es grande (\(np\ge5\), \(nq\ge5\)).

\seccion{5. Nivel de confianza y valor crítico}
\begin{formula}
\[P\bigl(-z_{\alpha/2}\le Z\le z_{\alpha/2}\bigr)=1-\alpha\]
\end{formula}
\begin{center}
\begin{tabular}{c|c|c}
\(1-\alpha\) & \(\alpha/2\) & \(z_{\alpha/2}\)\\ \hline
90\,\% & 0,05 & 1,645\\
95\,\% & 0,025 & 1,96\\
99\,\% & 0,005 & 2,576
\end{tabular}
\end{center}
Se obtiene con \(\Phi(z_{\alpha/2})=1-\frac\alpha2\).

\seccion{6. Intervalo característico}
Intervalo centrado en \(\mu\) que contiene la media muestral con probabilidad \(1-\alpha\):
\begin{formula}
\[\Bigl[\mu-z_{\alpha/2}\frac{\sigma}{\sqrt n};\ \ \mu+z_{\alpha/2}\frac{\sigma}{\sqrt n}\Bigr]\]
\end{formula}
Se usa cuando se \textbf{conoce} \(\mu\) y se pregunta por las muestras. Para un individuo, el radio es \(z_{\alpha/2}\sigma\).

\seccion{7. Intervalo de confianza para la media}
Se \textbf{desconoce} \(\mu\) y se estima a partir de \(\bar x\) (\(\sigma\) conocida):
\begin{formula}
\[IC_{1-\alpha}(\mu)=\bigl[\bar x-E;\ \ \bar x+E\bigr],\qquad E=z_{\alpha/2}\frac{\sigma}{\sqrt n}\]
\end{formula}
\begin{ejemplo}
\(\bar x=172\) cm, \(\sigma=10\), \(n=100\), confianza 95\,\%: \(E=1{,}96\cdot\frac{10}{10}=1{,}96\).
\[IC_{0{,}95}(\mu)=[170{,}04;\ 173{,}96]\]
\end{ejemplo}
\begin{ojo}
El 95\,\% se refiere al \textbf{método}: el 95\,\% de los intervalos construidos así contienen \(\mu\). No significa que \(\mu\) tenga probabilidad 0,95 de estar en este intervalo concreto.
\end{ojo}

\seccion{8. Intervalo para una proporción}
\begin{formula}
\[IC_{1-\alpha}(p)\approx\Bigl[\hat p-z_{\alpha/2}\sqrt{\tfrac{\hat p(1-\hat p)}n};\ \ \hat p+z_{\alpha/2}\sqrt{\tfrac{\hat p(1-\hat p)}n}\Bigr]\]
\end{formula}
\begin{ejemplo}
120 de 400 encuestados opinan que sí: \(\hat p=0{,}3\).
\(E=1{,}96\sqrt{\frac{0{,}3\cdot0{,}7}{400}}\approx0{,}0449\Rightarrow IC_{0{,}95}(p)\approx[0{,}2551;\ 0{,}3449]\).
\end{ejemplo}

\seccion{9. Error y tamaño de la muestra}
\begin{formula}
Media: \(\displaystyle n\ge\Bigl(\frac{z_{\alpha/2}\,\sigma}{E_0}\Bigr)^2\)\\[3pt]
Proporción: \(\displaystyle n\ge\frac{z_{\alpha/2}^2\,p_0(1-p_0)}{E_0^2}\); sin información, \(p_0=0{,}5\): \(\displaystyle n\ge\frac{z_{\alpha/2}^2}{4E_0^2}\)
\end{formula}
Se redondea siempre \textbf{hacia arriba}.
\begin{ejemplo}
\(\sigma=10\), \(E_0=2\), 95\,\%: \(n\ge\bigl(\frac{1{,}96\cdot10}2\bigr)^2=96{,}04\Rightarrow n=97\).
Proporción, \(E_0=0{,}03\), 95\,\%, sin información: \(n\ge\frac{1{,}96^2}{4\cdot0{,}03^2}\approx1067{,}1\Rightarrow n=1068\).
\end{ejemplo}
\sub{Relaciones.} \(E\) aumenta con el nivel de confianza y con \(\sigma\), y disminuye con \(n\). Como depende de \(\sqrt n\), para reducir el error a la mitad hay que multiplicar \(n\) por 4.

\seccion{10. Ficha técnica de una encuesta}
\begin{tabular}{@{}p{2.6cm}p{5.2cm}@{}}
Universo & Población estudiada\\
Ámbito & Lugar geográfico\\
Tamaño muestral & \(n\) (y entrevistas válidas)\\
Procedimiento & Técnica de muestreo y afijación\\
Error muestral & \(\pm E\), con su nivel de confianza\\
Nivel de confianza & 95\,\% habitual; supuesto \(p=q=0{,}5\)\\
Fecha y entidad & Trabajo de campo y responsable
\end{tabular}

\begin{ejemplo}
Ficha: \(n=1068\), muestreo aleatorio, 95\,\%, \(p=q=0{,}5\) \(\Rightarrow\) error máximo \(\pm3\)\,\%. Si el muestreo no es probabilístico, ese error no está justificado.
\end{ejemplo}

\seccion{11. Intervalos en la campana}
\begin{center}
\begin{tikzpicture}
\begin{axis}[mini,width=6.8cm,height=3.3cm,xmin=-3.6,xmax=3.6,ymin=0,ymax=0.45,xtick={-1.96,0,1.96},xticklabels={\(-z_{\alpha/2}\),\(0\),\(z_{\alpha/2}\)},tick label style={font=\tiny},xlabel={},ylabel={},axis y line=none]
\addplot[name path=c,domain=-3.6:3.6]{exp(-x^2/2)/sqrt(2*pi)};
\path[name path=ax] (axis cs:-3.6,0) -- (axis cs:3.6,0);
\addplot[fill=acento!30,draw=none] fill between[of=c and ax,soft clip={domain=-1.96:1.96}];
\addplot[fill=aviso!25,draw=none] fill between[of=c and ax,soft clip={domain=-3.6:-1.96}];
\addplot[fill=aviso!25,draw=none] fill between[of=c and ax,soft clip={domain=1.96:3.6}];
\node[font=\tiny] at (axis cs:0,0.15) {\(1-\alpha\)};
\node[font=\tiny] at (axis cs:-2.7,0.06) {\(\alpha/2\)};
\node[font=\tiny] at (axis cs:2.7,0.06) {\(\alpha/2\)};
\end{axis}
\end{tikzpicture}
\end{center}
\begin{ejemplo}
\textbf{Característico.} \(\mu=170\), \(\sigma=10\), \(n=25\), 95\,\%: el 95\,\% de las medias muestrales cae en \(170\pm1{,}96\cdot2=[166{,}08;\ 173{,}92]\).
\end{ejemplo}
\begin{ejemplo}
\textbf{Efecto de la confianza.} Con \(\bar x=172\), \(\sigma=10\), \(n=100\): al 90\,\%, \(E=1{,}645\); al 95\,\%, \(E=1{,}96\); al 99\,\%, \(E=2{,}576\). Más confianza, intervalo más ancho.
\end{ejemplo}

\seccion{12. Problema inverso}
Dado el intervalo \([170{,}04;\ 173{,}96]\) con \(\sigma=10\) y \(n=100\):
\begin{itemize}
  \item \(\bar x\) es el centro: \(\frac{170{,}04+173{,}96}2=172\).
  \item \(E\) es la mitad de la amplitud: \(1{,}96\).
  \item \(z_{\alpha/2}=\frac{E\sqrt n}{\sigma}=1{,}96\Rightarrow\) confianza del 95\,\%.
\end{itemize}
\begin{ejemplo}
\textbf{Tamaño con información previa.} Una encuesta anterior dio \(p_0=0{,}3\). Para \(E_0=0{,}03\) al 95\,\%: \(n\ge\frac{1{,}96^2\cdot0{,}3\cdot0{,}7}{0{,}03^2}\approx896{,}4\Rightarrow n=897\), menos que los 1068 del caso conservador.
\end{ejemplo}

\seccion{13. Traducir el enunciado}
\begin{tabular}{@{}p{3.6cm}p{4.2cm}@{}}
\textbf{Enunciado} & \textbf{Qué hay que hacer}\\ \hline
\guillemotleft{}Probabilidad de que la media de una muestra\ldots\guillemotright{} & Distribución de \(\bar X\)\\
\guillemotleft{}Entre qué valores estará la media muestral\ldots\guillemotright{} & Intervalo característico\\
\guillemotleft{}Estime con una confianza del\ldots\guillemotright{} & Intervalo de confianza\\
\guillemotleft{}\textquestiondown{}Qué tamaño debe tener la muestra\ldots?\guillemotright{} & \(n\) mínimo, hacia arriba\\
\guillemotleft{}Error máximo cometido\guillemotright{} & \(E=z_{\alpha/2}\sigma/\sqrt n\)
\end{tabular}

\seccion{14. Estrategia}
\begin{enumerate}
  \item \textquestiondown{}Se conoce el parámetro y se pregunta por las muestras? \(\to\) distribución muestral o intervalo característico.
  \item \textquestiondown{}Se desconoce y se estima? \(\to\) intervalo de confianza.
  \item \textquestiondown{}Se fija el error? \(\to\) tamaño de muestra, redondeando hacia arriba.
  \item Interpreta el resultado en el contexto, con unidades.
\end{enumerate}

\begin{ojo}
\begin{itemize}
  \item Usar \(\sigma\) en lugar de \(\frac\sigma{\sqrt n}\) para la media muestral.
  \item Redondear \(n\) hacia abajo.
  \item Confundir el intervalo característico (centrado en \(\mu\)) con el de confianza (centrado en \(\bar x\)).
  \item Elegir \(z=1{,}645\) para el 95\,\% bilateral.
\end{itemize}
\end{ojo}

\end{multicols}
\end{document}
